\scanpage{29}\mspage{28}

\deleted{\unsure{Alors}}

Ainsi

\uline{Proposition} On a construit un Inf de $F$ et $G$ pour $\ll$, on l'appelle la ``figure intersection'' $F.G$.

Attention, ce \uline{n'est pas} \inserted{en général} la figure $F\cap G$, \deleted{\unsure{quand|inter.|intersection}} en se rappelant $F,\deleted{F}G\subset\unsure{\mathfrak P|\mathscr P|\mathbb P}(\unsure{L,|L_1|L})$, \deleted{$G\subset\mathfrak P(L)$}

Nous n'allons considérer cette $F\cap G$ que dans le cas où $F$ et $G$ compatibles, mais alors les deux coïncident~!

Donc la bonne notion d'intersection est le $F.G$. En fait, ça marche pour une famille quelconque de figures, et
\[
 \left|\bigwedge_{i\in I} F_i\right|=\bigcap_{i\in I}|F_i|\quad .
\]

Cette propriété d'existence d'un Inf pour la relation $\ll$ \unsure{l'apparente|s'apparente} \deleted{\unsure{la propriété|la présentation}} à une propriété de ``\uline{modération}''.

\deleted{\uline{Propriété} Si $F'\preceq F$ \deleted{\unsure{, et des}}, alors $F'.G\preceq F.G$ (Immédiat sur supports)~: Subdivision \uline{induite} sur un raffinement.}
